\scanpage{75}
\textcircled{3} \deleted{$d'=$ \unsure{tous les|toutes les}} $d'$ \deleted{formé de diagrammes finis} \inserted{quelconque}, $d''=\varnothing$, \uline{pas de conditions d'exactitude}.

\deleted{\unsure{Pb 1} On suppose que} Soit $d'_0$ l'ens de tous les diagrammes finis, $d'_1$ l'ens de tous les petits diagrammes, et $\Delta\supset\Delta_0\supset\Delta_1$ les inclusions correspondantes.

\uline{Pb 1} $\Delta$-Enveloppe \inserted{$\rho_{\Delta}(I)$} d'une catégorie $I$. Soit $\overline I$ \deleted{$\unsure{\rho(I)}$ \illegible{1}} $\simeq\widehat{I^\circ}{}^\circ$ la $\Delta_1$ enveloppe, et soit $\widetilde I$ la \deleted{sous}-catégorie pl. de $\overline I$ formée des $\varprojlim$ de diagr. $d'$ d'objets de $I$ [NB $I\longrightarrow\overline I$ pl fid !].

Je dis que cela marche, \uline{pourvu} qu'on sache que $\widetilde I$ est stable par $d'$ (c'est le cas si $d'=$ ens. des diagrammes finis discrets, ou $d'=$ ens. de tous les diagrammes finis \inserted{(!)} \ldots) sinon on doit itérer et prendre \deleted{la lim} pour $\widetilde I$ la plus petite sous catégorie \inserted{str} pleine de $\overline I$ contenant $I$ et stable par $\varprojlim$ de type $d'$. En effet, on a pour $C\in\mathrm{Ob}\,\Delta$
% slanted
\margin{Pb 1 et Pb $1'$ Cela marche pour \uline{tout} ens. $d'$ de petits diagrammes (pas néc. finis)}
\[
\deleted{\underline{\mathrm{Hom}}_{\deleted{\unsure{\Delta|\mathrm{Cat}}}\mathrm{Cat}}(I,\unsure{\mathscr C|C})\longrightarrow\underline{\mathrm{Hom}}_{\Delta_1}(\overline I,\overline C),}
\]
\deleted{\deleted{et} \inserted{foncteur (pl. fidèle dont)} l'image ess. est formée de(s) foncteurs $\overline I\xrightarrow{\overline\varphi}\overline C$ tels que $\overline\varphi(I)\subset\deleted{\unsure{\mathrm{Im}}} C$ (ess.) Mais \deleted{alors} $\overline\varphi^{-1}(C)$}

Soit $\rho(C)$ la $\Delta$-$\Delta_1$ enveloppe de $C$ i.e. la $\Delta_1$-catégorie avec $\Delta$-foncteur $C\xrightarrow{i}\rho(C)$ \uline{universel} [c'est une catégorie \deleted{\unsure{quot.|pleine}} de fraction de $\overline C$]. \deleted{\unsure{Alors supposons que}} \inserted{Admettons que} \deleted{$I\longrightarrow C$ se prolonge} ce foncteur est pleinement fidèle (Pb auxiliaire) \deleted{En effet, \unsure{c'est} les foncteurs} Alors les foncteurs $\varphi:I\longrightarrow C$ \deleted{définist, \unsure{via i$\varphi$|à $\varphi$}} \unsure{``s.a''|``s.a. à''} les foncteurs $\psi=i\varphi:\widetilde I\longrightarrow\rho(C)$ qui envoient $I$ dans l'image essentielle de $C$ par $i$, ou encore les foncteurs $\overline\psi:\overline I\longrightarrow\rho(C)$. tels que $\overline\psi(I)$ soit contenue dans l'image essentielle de $\unsure{I|i|R|C}$ ; de tels $\overline\psi$ envoient $\widetilde I$ dans $C$ d'où $\underline{\mathrm{Hom}}_{\mathrm{Cat}}(I,C)\longrightarrow\underline{\mathrm{Hom}}_{\Delta}(\widetilde I,C)$, et on a un foncteur ``restriction'' en sens inverse. Il est clair que les composés dans les deux sens sont l'identité, OK.

\uline{Pb $1'$} \uline{Ou} Pb \uline{auxiliaire} Si $C\in\mathrm{Ob}\,\Delta$, montrer que $C\longrightarrow\rho(C)$ \uline{est} \uline{pl.} \uline{fidèle}. Je dis qu'on a
\[
\rho(C)=\underline{\mathrm{Hom}}_\Delta(C,\mathrm{Ens})^\circ=T_C(\mathrm{Ens})^\circ,
\]
% archive sheet number: 74
